\scanpage{23}
\mspage{- 218 -}\hfill n\textsuperscript{o} 391

(VII) La demi-somme des racines positives est la somme des poids fondamentaux, (\S~6, n\textsuperscript{o}~9, prop. 8), et vaut donc :

\[
\begin{aligned}
\rho&=\varepsilon_2+2\varepsilon_3+3\varepsilon_4+4\varepsilon_5+4(\varepsilon_8-\varepsilon_7-\varepsilon_6)=\\
&=8\alpha_1+11\alpha_2+15\alpha_3+21\alpha_4+15\alpha_5+8\alpha_6.
\end{aligned}
\]

(VIII) De même qu'au n\textsuperscript{o}~11, le réseau des vocables $\unsure{\Gamma_o|\Gamma_0}(R'')$ est l'associé \uline{dans} $V''$ du réseau $\Gamma_2^*$ de $V''{}^*$ \uline{projection orthogonale} dans $V''{}^*$ de $\Gamma_o(R^*)$, (lorsque $V''{}^*$ est identifié au sous-espace de $V^*$ orthogonal à $\varepsilon_7^*+\varepsilon_8^*$ et $\varepsilon_7^*-\varepsilon_6^*$). Le groupe $\Gamma(R'')/\Gamma_o(R'')$ est donc \uline{isomorphe au groupe} $\Gamma_2^*/\Gamma_o(R''{}^*)$, qu'il reste à calculer. Or, les projections orthogonales de $\varepsilon_8^*$, $-\varepsilon_7^*$ et $-\varepsilon_6^*$ sur $V''{}^*$ sont toutes trois égales $\frac{1}{3}(\varepsilon_8^*-\varepsilon_7^*-\varepsilon_6^*)$ et les autres $\varepsilon_i^*$ sont dans $V''{}^*$ ; on en déduit aisément que $\Gamma_2^*$ est engendré par $\Gamma_o(R''{}^*)$ et par l'unique vecteur $\mu^*=\frac{1}{3}(\varepsilon_8^*-\varepsilon_7^*-\varepsilon_6^*)$ ; mais comme $3\mu^*\in\Gamma_o(R''{}^*)$, on voit que $\Gamma_2^*/\Gamma_o(R''{}^*)$ est isomorphe à $Z/3Z$ et par suite $\Gamma(R'')/\Gamma_o(R'')$ est \uline{isomorphe à} $Z/3Z$, et l'indice de connexion $f=3$.

(IX) et (X). Le calcul de la plus grande racine et de l'indice de connexion permet, grâce au \S~6, n\textsuperscript{o}~10, prop. 23, de calculer l'ordre du groupe de Weyl $W(R'')$, qui est égal à :

\[
6!1.2.2.3.2.1.3=2^7.3^4.5=51.840\unsure{.|:}
\]

La suite des exposants a 6 termes compris entre 1 et 11, et comprend les termes 1, 5, 7, 11 qui sont étrangers à 12, (\S~8, n\textsuperscript{o}~5, prop. 6). Les deux derniers termes de la suite $m$, $m'$ sont des entiers tels que :

\[
\begin{aligned}
m+m'&=12\\
(m+1)(m'+1)(1+1)(5+1)(7+1)(12+1)&=2^7.3^4.5
\end{aligned}
\]

% archive sheet number: 22
